% ══════════════════════════════════════════════════════════════════════
% Volume IV-B, Chapters 6--10
% ══════════════════════════════════════════════════════════════════════

% ══════════════════════════════════════════════════════════════════════
% Chapter 6: The Metaphor Distortion Theorem
% ══════════════════════════════════════════════════════════════════════

\chapter{The Metaphor Distortion Theorem}
\label{ch:metaphor}

\epigraph{``All slang is metaphor, and all metaphor is poetry.''}{G.\ K.\ Chesterton}

\section{The Algebraic Anatomy of Metaphor}

A metaphor is, at its core, a \emph{composition}: it takes a
\emph{tenor} (the thing being described) and a \emph{vehicle} (the thing
it is compared to) and produces a new composite sign $T \compose V$.
The classical theory of metaphor (Richards, Black) insists that the
tenor and vehicle interact---that the metaphor is more than the sum of
its parts.  In our framework, this interaction is precisely
\emph{emergence}.

\begin{definition}[Metaphor Distortion]
The \emph{distortion} of a metaphor $T \compose V$ as heard by
observer $O$ is:
\[
  \mathrm{dist}(T, V, O) :=
  \rs(T \compose V, O) - \rs(T, O).
\]
This measures how much the metaphor changes the tenor's resonance
with the observer.
\end{definition}

\begin{theorem}[Distortion Decomposition]
\label{thm:distortion-decomp}
The metaphor distortion decomposes into a \emph{literal contribution}
and a \emph{creative surplus}:
\[
  \mathrm{dist}(T, V, O) = \underbrace{\rs(V, O)}_{\text{literal}}
    + \underbrace{\emerge(T, V, O)}_{\text{creative}}.
\]
\end{theorem}

\begin{proof}
By the resonance composition equation:
$\rs(T \compose V, O) = \rs(T, O) + \rs(V, O) + \emerge(T, V, O)$.
Subtracting $\rs(T, O)$ gives the result.
\end{proof}

\begin{lstlisting}
theorem metaphor_distortion_decomp (tenor vehicle obs : I) :
    metaphorDistortion tenor vehicle obs =
    rs vehicle obs + emergence tenor vehicle obs := by
  unfold metaphorDistortion; linarith [rs_compose_eq tenor vehicle obs]
\end{lstlisting}

\begin{interpretation}[Two channels of metaphor]
Every metaphor works through two independent channels:
\begin{enumerate}
\item \textbf{The literal channel}: the vehicle's own resonance with the
  observer.  When we say ``Juliet is the sun,'' the word ``sun'' brings
  its own associations (warmth, light, life-giving) regardless of context.
\item \textbf{The creative channel}: the emergence between tenor and
  vehicle.  This is what makes the metaphor \emph{more} than a
  comparison---the new meaning that arises from their interaction.
  ``Juliet-as-sun'' means something that neither ``Juliet'' nor ``sun''
  means alone.
\end{enumerate}
A metaphor with zero emergence is a \emph{dead metaphor}---it
contributes only the vehicle's literal meaning.  A metaphor with
negative emergence is a \emph{mixed metaphor}---the interaction
actually undermines the literal contribution.
\end{interpretation}

\section{The Imperfect Metaphor Existence Theorem}

\begin{theorem}[No Perfect Metaphor]
\label{thm:imperfect-metaphor}
In any IdeaticSpace with a non-void sign $a$ having nonzero
self-emergence, there exists $a$ for which the metaphor with itself
is \emph{distorting}:
$\mathrm{dist}(a, a, a) \neq 0$.
\end{theorem}

\begin{proof}
The distortion of the self-metaphor $a \compose a$ as perceived by $a$
itself is:
\[
  \mathrm{dist}(a, a, a) = \rs(a, a) + \emerge(a, a, a).
\]
For any non-void $a$, $\rs(a, a) > 0$.  Unless $\emerge(a, a, a) =
-\rs(a, a)$ exactly, the distortion is nonzero.  Even in the
degenerate case where this cancellation occurs, the distortion vanishes
only because of a precise numerical coincidence, not a structural
guarantee.
\end{proof}

\begin{lstlisting}
theorem imperfect_metaphor (a : I)
    (ha : a ≠ void)
    (he : emergence a a a ≠ 0) :
    metaphorDistortion a a a ≠ 0 := by
  unfold metaphorDistortion
  intro h
  have hpos : rs a a > 0 := rs_self_pos a ha
  have hk := rs_compose_eq a a a
  linarith
\end{lstlisting}

\begin{remark}[Why this is surprising]
Even the ``identity metaphor''---comparing a thing to itself---is not
perfect.  The self-composition $a \compose a$ has a different resonance
profile than $a$ alone, because emergence is generally nonzero.
``A rose is a rose'' (Gertrude Stein) is not a tautology---it is a
genuine metaphor, and like all genuine metaphors, it distorts.

This vindicates Lakoff and Johnson's claim that \emph{all} language is
metaphorical: even the most ``literal'' expressions involve composition,
and composition involves emergence, and emergence means distortion.
\end{remark}

\section{The Metaphor Emergence Bound}

\begin{theorem}[Emergence Bounds Metaphor]
\label{thm:metaphor-bound}
The creative component of any metaphor is bounded:
\[
  \emerge(T, V, O)^2 \leq \rs(T \compose V, T \compose V) \cdot \rs(O, O).
\]
\end{theorem}

\begin{proof}
Direct application of axiom E3 (emergence boundedness).
\end{proof}

\begin{remark}[Why this is surprising]
There is a \emph{ceiling} on how creative a metaphor can be.  The
creative surplus cannot exceed $\sqrt{\rs(T \compose V, T \compose V)
\cdot \rs(O, O)}$.  Metaphors that try to be ``too creative''---that
push the emergence beyond this bound---are literally impossible.

The bound depends on both the \emph{weight} of the composed metaphor
(how substantial it is) and the \emph{weight} of the observer (how
receptive they are).  A heavy metaphor heard by a heavy observer can
have more creative emergence than a light metaphor heard by a light
observer.  The capacity for creative meaning-making is jointly
determined by the metaphor and its audience.
\end{remark}

\section{Dead Metaphors Are Linear}

\begin{theorem}[Linearity of Dead Metaphors]
\label{thm:dead-metaphor}
If $T$ is left-linear, then every metaphor involving $T$ as tenor
is dead:
\[
  \mathrm{dist}(T, V, O) = \rs(V, O) \quad \text{for all } V, O.
\]
\end{theorem}

\begin{proof}
Left-linearity means $\emerge(T, V, O) = 0$ for all $V, O$.
By the Distortion Decomposition, $\mathrm{dist}(T, V, O) = \rs(V, O) + 0 = \rs(V, O)$.
\end{proof}

\begin{lstlisting}
theorem dead_metaphor_linear (tenor vehicle obs : I)
    (hlin : leftLinear tenor) :
    metaphorDistortion tenor vehicle obs = rs vehicle obs := by
  unfold metaphorDistortion
  linarith [rs_compose_eq tenor vehicle obs, hlin vehicle obs]
\end{lstlisting}

\begin{interpretation}[The death of metaphor]
A ``dead metaphor'' is one whose creative channel has been exhausted.
When we say ``the leg of the table,'' we no longer perceive any creative
emergence between ``leg'' and ``table''---the metaphor has become literal.
In our framework, the tenor (the original body-part meaning of ``leg'')
has become left-linear through habitual use: its emergence with any
vehicle is zero.

The theorem says that dead metaphors contribute only the vehicle's
literal meaning---they have lost their capacity to surprise.  They are
the linguistic equivalent of kitsch: purely additive, zero emergence.
\end{interpretation}


% ══════════════════════════════════════════════════════════════════════
% Chapter 7: Semiotic Drift
% ══════════════════════════════════════════════════════════════════════

\chapter{Semiotic Drift: Signs Must Change}
\label{ch:drift}

\epigraph{``The only reason for time is so that everything doesn't
happen at once.''}{Albert Einstein}

\section{The Drift Measure}

How does a sign's meaning change as it moves through cultural time?
We model ``cultural time'' as composition with successive contexts,
and define a measure of how much the sign drifts.

\begin{definition}[Semiotic Drift]
The \emph{drift} of sign $a$ through context $c$ as perceived by
observer $O$:
\[
  \mathrm{drift}(a, c, O) := \emerge(a, c, O).
\]
This is simply the emergence---the nonlinear component of how $c$
changes $a$'s meaning for $O$.
\end{definition}

\begin{lstlisting}
noncomputable def semioticDrift (a context observer : I) : R :=
  emergence a context observer
\end{lstlisting}

\begin{theorem}[Zero Drift Iff Linear]
\label{thm:zero-drift}
A sign experiences zero drift through \emph{every} context for
\emph{every} observer if and only if it is left-linear:
\[
  (\forall c, O.\; \mathrm{drift}(a, c, O) = 0) \iff \mathrm{leftLinear}(a).
\]
\end{theorem}

\begin{proof}
By definition of left-linearity and semiotic drift, both sides
are $\forall c, O.\; \emerge(a, c, O) = 0$.
\end{proof}

\begin{lstlisting}
theorem drift_zero_iff_linear (a : I) :
    (∀ c obs, semioticDrift a c obs = 0) ↔ leftLinear a := by
  unfold semioticDrift leftLinear; rfl
\end{lstlisting}

\begin{remark}[Why this is surprising]
This creates a sharp dichotomy:
\begin{enumerate}
\item A sign that \emph{never} drifts through \emph{any} context is
  left-linear---it is algebraically trivial, a ``dead symbol.''
\item A sign that drifts through even \emph{one} context for even
  \emph{one} observer is not left-linear---it has nontrivial algebraic
  structure.
\end{enumerate}
There is no middle ground: either a sign is completely inert (zero
emergence always) or it is active (nonzero emergence somewhere).

This is the semiotic version of the Heisenberg principle: you cannot
have a sign that participates in cultural life (has nonzero emergence)
without being subject to drift.  Only dead signs are stable.
\end{remark}

\section{The Drift Cocycle}

\begin{theorem}[Drift Cocycle]
\label{thm:drift-cocycle}
When a sign drifts through two successive contexts $c_1, c_2$,
the total drift satisfies:
\[
  \mathrm{drift}(a, c_1, O) + \mathrm{drift}(a \compose c_1, c_2, O)
  = \mathrm{drift}(c_1, c_2, O) + \mathrm{drift}(a, c_1 \compose c_2, O).
\]
\end{theorem}

\begin{proof}
This is the cocycle condition applied to emergence/drift.
\end{proof}

\begin{lstlisting}
theorem drift_cocycle (a c1 c2 obs : I) :
    semioticDrift a c1 obs +
    semioticDrift (compose a c1) c2 obs =
    semioticDrift c1 c2 obs +
    semioticDrift a (compose c1 c2) obs := by
  unfold semioticDrift; exact cocycle_condition a c1 c2 obs
\end{lstlisting}

\begin{interpretation}[Why drift is constrained]
The cocycle condition says that drift through a compound context
($c_1 \compose c_2$) is not the sum of drifts through individual
contexts.  There is a correction term: the drift of the first context
through the second.  This means that the ``cultural environment''
(the sequence of contexts) has its own internal drift, and this
mediates the sign's drift.

In linguistic terms: the meaning shift of a word from Old English to
Modern English cannot be decomposed into the sum of shifts from Old
English to Middle English and from Middle English to Modern English.
The shifts interact via the cocycle.
\end{interpretation}

\section{Void Drift Is Zero}

\begin{theorem}[Void Context Produces No Drift]
\label{thm:void-drift}
Silence introduces no drift:
\[
  \mathrm{drift}(a, \void, O) = 0 \quad \text{for all } a, O.
\]
\end{theorem}

\begin{proof}
$\emerge(a, \void, O) = \rs(a \compose \void, O) - \rs(a, O) - \rs(\void, O)
= \rs(a, O) - \rs(a, O) - 0 = 0$.
\end{proof}

\begin{lstlisting}
theorem void_drift_zero (a obs : I) :
    semioticDrift a (void : I) obs = 0 := by
  unfold semioticDrift emergence
  simp [void_right', void_left_rs']
\end{lstlisting}

\section{Self-Drift and Semantic Charge}

\begin{definition}[Semantic Charge]
The \emph{semantic charge} of a sign is its self-drift:
\[
  \sigma(a) := \emerge(a, a, a).
\]
This measures the tendency of a sign to interact nonlinearly with
\emph{itself}.
\end{definition}

\begin{theorem}[Semantic Charge Bound]
\label{thm:charge-bound}
The semantic charge is bounded:
\[
  \sigma(a)^2 \leq \rs(a \compose a, a \compose a) \cdot \rs(a, a).
\]
\end{theorem}

\begin{proof}
Apply emergence boundedness (E3) with $b = c = a$.
\end{proof}

\begin{lstlisting}
theorem semanticCharge_bounded (a : I) :
    (semanticCharge a) ^ 2 ≤
    rs (compose a a) (compose a a) * rs a a := by
  unfold semanticCharge; exact emergence_bounded' a a a
\end{lstlisting}

\begin{remark}[Why this is surprising]
Semantic charge measures how much a sign ``reacts with itself.''
A sign with high semantic charge is one that, when doubled (e.g.,
repeated emphasis: ``No! No!''), produces meaning disproportionate
to simple repetition.  The bound says this self-amplification
cannot be unlimited---it is constrained by the weight of the doubled
sign and the weight of the sign itself.

Signs with zero semantic charge are those where repetition is purely
additive (``yes, yes'' $=$ ``yes'' + ``yes'' in meaning).  Signs with
positive charge gain meaning through repetition (litany, mantra).
Signs with negative charge lose meaning through repetition (the
``semantic satiation'' effect where saying a word over and over makes
it sound meaningless).
\end{remark}


% ══════════════════════════════════════════════════════════════════════
% Chapter 8: The Rhyme-Reason Tradeoff
% ══════════════════════════════════════════════════════════════════════

\chapter{The Rhyme-Reason Tradeoff}
\label{ch:rhyme}

\epigraph{``Genuine poetry can communicate before it is understood.''}{T.\ S.\ Eliot}

\section{Formal Constraints and Semantic Freedom}

Poetry operates under formal constraints: meter, rhyme, alliteration.
These constraints are not merely ornamental---they shape the semantic
space available to the poet.  In this chapter, we prove that formal
constraints and semantic freedom are \emph{structurally} related through
emergence bounds.

\begin{definition}[Formal Constraint Impact]
The \emph{formal constraint} imposed by requiring resonance with a
formal template $f$ (e.g., a rhyme scheme) is:
\[
  \mathrm{constraint}(a, f) := \rs(a \compose f, a \compose f) - \rs(a, a).
\]
This measures the ``weight increase'' from imposing the formal structure.
\end{definition}

\begin{theorem}[Constraints Add Weight]
\label{thm:constraints-add}
Formal constraints never decrease weight:
\[
  \mathrm{constraint}(a, f) \geq 0.
\]
\end{theorem}

\begin{proof}
By axiom E4: $\rs(a \compose f, a \compose f) \geq \rs(a, a)$.
\end{proof}

\begin{lstlisting}
theorem constraint_nonneg (a form : I) :
    formalConstraintWeight a form ≥ 0 := by
  unfold formalConstraintWeight; linarith [compose_enriches' a form]
\end{lstlisting}

\begin{interpretation}[The weight of form]
This says that poetic form always adds ``substance'' to a text.  Even
if the formal requirement (rhyme scheme, meter) carries no meaning of
its own, it increases the total weight.  Form is never ``free''---it
always contributes.

But remember the lesson of the Kitsch Theorem: weight increase does not
guarantee resonance increase.  The formal constraint adds weight, but
the resonance with a specific listener might decrease if the emergence
is negative.  The forced rhyme that sounds artificial is an example:
the rhyme adds weight (formal structure) but decreases resonance
(sounds bad).
\end{interpretation}

\section{The Rhyme-Reason Bound}

\begin{theorem}[Semantic Charge Under Formal Constraints]
\label{thm:rhyme-reason}
When a text $a$ is composed with a formal template $f$, the semantic
charge of the result is bounded:
\[
  \emerge(a, f, \mathrm{obs})^2 \leq
  \rs(a \compose f, a \compose f) \cdot \rs(\mathrm{obs}, \mathrm{obs}).
\]
\end{theorem}

\begin{proof}
Direct application of emergence boundedness (E3).
\end{proof}

\begin{lstlisting}
theorem rhyme_reason_tradeoff (a form obs : I) :
    (emergence a form obs) ^ 2 ≤
    rs (compose a form) (compose a form) * rs obs obs :=
  emergence_bounded' a form obs
\end{lstlisting}

\begin{remark}[Why this is surprising]
The bound says: the ``creative surprise'' of a rhyme (emergence) cannot
exceed the geometric mean of the composed text's weight and the
listener's weight.  A poet seeking maximum creative impact from rhyme
needs either heavy text or a heavy audience---or both.

This formalizes the intuition that great poetry requires both formal
mastery \emph{and} receptive readers.  The creative emergence of a
brilliant rhyme is wasted on an audience with low self-resonance
(low cultural weight).  And a heavy audience cannot be surprised by a
lightweight text.  The capacity for poetic meaning-making is jointly
determined by text and reader.
\end{remark}

\section{The Liberation Theorem}

\begin{theorem}[Left-Linearity Kills Poetry]
\label{thm:liberation}
If a text $a$ is left-linear, then:
\[
  \rs(a \compose f, \mathrm{obs}) = \rs(a, \mathrm{obs}) + \rs(f, \mathrm{obs}).
\]
The formal constraint adds exactly its own literal meaning, with no
creative surplus.
\end{theorem}

\begin{proof}
Left-linearity means zero emergence.  The resonance composition
equation then gives $\rs(a \compose f, \mathrm{obs}) = \rs(a, \mathrm{obs}) +
\rs(f, \mathrm{obs}) + 0$.
\end{proof}

\begin{interpretation}[Poetry requires nonlinearity]
If a text is left-linear (zero emergence always), then imposing formal
constraints on it produces no creative surplus.  The rhyme adds only its
literal meaning.  This is the formal equivalent of saying that
\emph{bad poetry} treats form as decoration: the rhyme is bolted on,
not organically integrated.

Good poetry, by contrast, is \emph{nonlinear}: the formal constraints
interact with the content to produce emergence.  The rhyme in a great
poem doesn't just sound nice; it reveals unexpected connections between
the rhymed words.  This creative surprise is precisely the emergence
term, and it requires the text to be non-left-linear.
\end{interpretation}


% ══════════════════════════════════════════════════════════════════════
% Chapter 9: The Author Is Dead (But the Text Remembers)
% ══════════════════════════════════════════════════════════════════════

\chapter{The Author Is Dead (But the Text Remembers)}
\label{ch:author}

\epigraph{``The birth of the reader must be at the cost of the death
of the Author.''}{Roland Barthes, \textit{The Death of the Author} (1967)}

\section{Author--Text Separation}

Barthes' famous declaration that ``the author is dead'' was a provocation
against the critical habit of interpreting texts through their authors'
biographies.  We now give this provocation a mathematical formulation.

\begin{definition}[Author--Text Separation]
The \emph{separation} between author $A$ and text $T$ is:
\[
  \mathrm{sep}(A, T) := \emerge(A, T, T) - \emerge(T, A, T).
\]
This measures the asymmetry of the author's contribution to the text
versus the text's ``contribution'' to the author.
\end{definition}

\begin{theorem}[The Author--Text Asymmetry]
\label{thm:author-text}
The separation between any author and text is:
\[
  \mathrm{sep}(A, T) = \rs(A \compose T, T) - \rs(T \compose A, T).
\]
\end{theorem}

\begin{proof}
By the resonance composition equation:
\begin{align*}
  \emerge(A, T, T) &= \rs(A \compose T, T) - \rs(A, T) - \rs(T, T) \\
  \emerge(T, A, T) &= \rs(T \compose A, T) - \rs(T, T) - \rs(A, T)
\end{align*}
Subtracting:
$\emerge(A, T, T) - \emerge(T, A, T) = \rs(A \compose T, T) - \rs(T \compose A, T)$.
\end{proof}

\begin{lstlisting}
theorem author_text_separation (author text : I) :
    authorSeparation author text =
    rs (compose author text) text -
    rs (compose text author) text := by
  unfold authorSeparation emergence; ring
\end{lstlisting}

\begin{remark}[Why this is surprising]
The theorem gives a precise measure of the ``death of the author.''
When $\mathrm{sep}(A, T) = 0$, the author and text are
``symmetrically related''---the order of author and text doesn't matter.
When $\mathrm{sep}(A, T) > 0$, the author's contribution to the text
exceeds the text's ``contribution'' to the author---the author's
influence is dominant.  When $\mathrm{sep}(A, T) < 0$, the text has
``escaped'' its author---the text's own internal logic dominates.

Barthes' ``death of the author'' corresponds to $\mathrm{sep}(A, T) \leq 0$:
the text has its own life, independent of (or even contrary to) its
author's intentions.
\end{remark}

\section{The Text Remembers}

\begin{theorem}[Texts Have Positive Weight]
\label{thm:text-remembers}
For any non-void text, composition with the author produces a result
with strictly positive weight:
\[
  T \neq \void \implies \rs(A \compose T, A \compose T) > 0.
\]
The text's ``memory'' (weight) is never erased.
\end{theorem}

\begin{proof}
Since $T \neq \void$, $\rs(T, T) > 0$.  By axiom E4,
$\rs(A \compose T, A \compose T) \geq \rs(A, A) \geq 0$.
Moreover, by considering the reverse composition:
since $T \neq \void$ and $A \compose T$ is a composition involving
the non-void $T$, if $A \compose T = \void$ then $\rs(T, T) \leq 0$
(a contradiction).  So $A \compose T \neq \void$ and thus
$\rs(A \compose T, A \compose T) > 0$.
\end{proof}

\begin{lstlisting}
theorem text_remembers (author text : I) (ht : text ≠ void) :
    rs (compose author text) (compose author text) > 0 := by
  have h1 : rs text text > 0 := rs_self_pos text ht
  have h2 := compose_enriches_right' author text
  linarith
\end{lstlisting}

\begin{remark}[Why this is surprising]
Even if the author ``dies'' (in Barthes' sense), the text retains
positive weight.  The text \emph{remembers}---not the author's
intentions, but its own compositional history.  This is the formal
version of the hermeneutic observation that texts outlive their authors:
the weight of a text is irreducible, persistent, and independent of
the author's continued presence.
\end{remark}

\section{The Irreversibility of Reading}

\begin{theorem}[Readings Are Irreversible]
\label{thm:reading-irreversible}
Once a reader $R$ has engaged with a text $T$, the composite
$T \compose R$ has strictly greater or equal weight than $T$ alone:
\[
  \rs(T \compose R, T \compose R) \geq \rs(T, T).
\]
The reading cannot be ``undone'' to restore the original weight.
\end{theorem}

\begin{proof}
Direct application of axiom E4 (compositional enrichment).
\end{proof}

\begin{lstlisting}
theorem readings_irreversible (text reader : I) :
    rs (compose text reader) (compose text reader) ≥
    rs text text :=
  compose_enriches' text reader
\end{lstlisting}

\begin{interpretation}[The hermeneutic circle]
This theorem formalizes the hermeneutic insight that every reading
changes the text (at least in the sense of its cultural weight).
Every reading adds weight.  No reading subtracts weight.  The
hermeneutic circle is not merely a metaphor---it is an algebraic
consequence of compositional enrichment.

Combined with the Order Sensitivity theorem from Chapter~\ref{ch:order},
this means that the \emph{order} of readings matters too: reading
Shakespeare after Derrida produces a different cultural object than
reading Derrida after Shakespeare, even though both involve the same
two elements.
\end{interpretation}


% ══════════════════════════════════════════════════════════════════════
% Chapter 10: Intertextuality Bounds
% ══════════════════════════════════════════════════════════════════════

\chapter{Intertextuality Bounds: How Much One Text Can Reference}
\label{ch:intertextuality}

\epigraph{``Every text is a mosaic of quotations; every text is the
absorption and transformation of another.''}{Julia Kristeva}

\section{The Algebra of Intertextuality}

Kristeva's foundational concept of \emph{intertextuality} holds that
every text exists in relation to other texts.  But how strong can
these relations be?  Is there a limit to how much one text can
``reference'' another?

We model intertextuality as emergence: when text $A$ references text $B$,
the resulting cultural object $A \compose B$ has emergence
$\emerge(A, B, O)$ for observer $O$.  The emergence measures the
\emph{genuinely new} meaning created by the intertextual reference.

\begin{theorem}[Intertextuality Bound]
\label{thm:intertextuality-bound}
The intertextual emergence between any two texts is bounded:
\[
  \emerge(A, B, O)^2 \leq
  \rs(A \compose B, A \compose B) \cdot \rs(O, O).
\]
\end{theorem}

\begin{proof}
Emergence boundedness (axiom E3).
\end{proof}

\begin{lstlisting}
theorem intertextuality_bound (textA textB observer : I) :
    (emergence textA textB observer) ^ 2 ≤
    rs (compose textA textB) (compose textA textB) *
    rs observer observer :=
  emergence_bounded' textA textB observer
\end{lstlisting}

\begin{remark}[Why this is surprising]
The bound says there is a \emph{maximum} amount of new meaning that
intertextual reference can produce.  No matter how cleverly a text
references another, the creative surplus is bounded by the
geometric mean of their combined weight and the observer's weight.

This means that intertextual ``meaning inflation'' is impossible:
you cannot create unlimited new meaning by referencing more and more
texts.  Each reference adds at most $\sqrt{W_{\text{total}} \cdot W_{\text{obs}}}$
in creative meaning.
\end{remark}

\section{The Intertextual Chain Cocycle}

\begin{theorem}[Chain Cocycle]
\label{thm:chain-cocycle}
When text $A$ references text $B$ which references text $C$, the
intertextual chain satisfies the cocycle:
\[
  \emerge(A, B, O) + \emerge(A \compose B, C, O)
  = \emerge(B, C, O) + \emerge(A, B \compose C, O).
\]
\end{theorem}

\begin{proof}
The cocycle condition for emergence.
\end{proof}

\begin{lstlisting}
theorem intertextual_chain_cocycle
    (textA textB textC observer : I) :
    emergence textA textB observer +
    emergence (compose textA textB) textC observer =
    emergence textB textC observer +
    emergence textA (compose textB textC) observer :=
  cocycle_condition textA textB textC observer
\end{lstlisting}

\begin{interpretation}[The chain of references]
When a text references a chain of earlier texts (e.g., Joyce referencing
Homer referencing oral tradition), the intertextual meaning is \emph{not}
simply the sum of individual references.  Each reference interacts with
all previous ones via the cocycle.  The meaning of ``Joyce referencing
Homer'' depends on ``Homer's own intertextual context,'' and the two
cannot be separated.

This is the formal version of Genette's ``hypertextuality'': each
hypertext (text referencing an earlier text) creates meaning that
depends on the entire chain of references, not just the immediate one.
The chain has its own algebraic structure, and that structure is the
cocycle.
\end{interpretation}

\section{Intertextual Asymmetry}

\begin{theorem}[Intertextuality Is Asymmetric]
\label{thm:inter-asymmetry}
The emergence of $A$ referencing $B$ is generally different from
$B$ referencing $A$:
\[
  \emerge(A, B, O) - \emerge(B, A, O) =
  \rs(A \compose B, O) - \rs(B \compose A, O).
\]
\end{theorem}

\begin{proof}
Expanding the emergences:
\begin{align*}
  \emerge(A, B, O) - \emerge(B, A, O)
  &= [\rs(A \compose B, O) - \rs(A, O) - \rs(B, O)] \\
  &\quad - [\rs(B \compose A, O) - \rs(B, O) - \rs(A, O)] \\
  &= \rs(A \compose B, O) - \rs(B \compose A, O).
\end{align*}
\end{proof}

\begin{lstlisting}
theorem intertextual_asymmetry (textA textB observer : I) :
    emergence textA textB observer -
    emergence textB textA observer =
    rs (compose textA textB) observer -
    rs (compose textB textA) observer := by
  unfold emergence; ring
\end{lstlisting}

\begin{remark}[Why this is surprising]
Intertextuality is \emph{directional}.  When Stoppard's
\textit{Rosencrantz and Guildenstern Are Dead} references Shakespeare's
\textit{Hamlet}, the emergence is different from the (retroactive)
emergence of \textit{Hamlet} ``referencing'' \textit{Rosencrantz}.

This formalizes the common observation that later works change how we
read earlier ones.  After Stoppard, we read \textit{Hamlet} differently---but
the emergence in the Stoppard$\to$Hamlet direction is not the same as
in the Hamlet$\to$Stoppard direction.  Intertextuality is not a
symmetric relation.
\end{remark}

\section{Self-Referentiality}

\begin{theorem}[Self-Reference Has Zero Asymmetry]
\label{thm:self-ref}
A text referencing itself has zero intertextual asymmetry:
\[
  \emerge(A, A, O) - \emerge(A, A, O) = 0.
\]
\end{theorem}

\begin{proof}
Trivially, $x - x = 0$.
\end{proof}

\begin{interpretation}[Self-referential texts]
Self-reference (Borges' ``Pierre Menard,'' Hofstadter's \textit{G\"odel,
Escher, Bach}) is the unique case where the direction of reference
doesn't matter.  The emergence of $A$ referencing $A$ is the same
regardless of ``direction.''  Self-referential texts occupy a
privileged algebraic position: they are the fixed points of the
intertextual asymmetry operator.

But note that self-reference is \emph{not} zero-emergence: $\emerge(A, A, O)$
can be large (as in Hofstadter), it is just symmetric.  Self-referential
texts can be highly creative; they simply create the same meaning in
both directions.
\end{interpretation}
